\documentclass[11pt]{article}
\usepackage[T1]{fontenc}
\usepackage{lmodern}
\usepackage[margin=1in]{geometry}
\usepackage{amsmath,amssymb,amsthm,mathtools}
\usepackage{enumitem}
\usepackage{microtype}
\usepackage{tikz}
\usepackage[colorlinks=true,linkcolor=blue!45!black,citecolor=green!35!black,urlcolor=blue!55!black]{hyperref}
\usepackage[all]{hypcap}
\hypersetup{
  pdftitle={A product counterexample to the simplex maximum for projection-body volume},
  pdfsubject={An elementary counterexample in dimension twenty}
}
\newtheorem{theorem}{Theorem}
\newtheorem{lemma}[theorem]{Lemma}
\newtheorem{proposition}[theorem]{Proposition}
\newtheorem{corollary}[theorem]{Corollary}
\newcommand{\R}{\mathbb R}
\DeclareMathOperator{\conv}{conv}
\DeclareMathOperator{\vol}{vol}
\DeclareMathOperator{\proj}{proj}
\newcommand{\unverified}[1]{\par\smallskip\noindent\fbox{\parbox{0.95\linewidth}{\textbf{Unverified step.} #1}}\par\smallskip}
\newcommand{\step}[1]{\smallskip\noindent\textbf{#1}}
\title{A product counterexample to the simplex maximum\\
for projection-body volume}
\author{Rewritten for readability from the OpenAI manuscript of the same title}
\date{September 24, 2026}
\raggedbottom

\begin{document}
\maketitle

\begin{abstract}
The product of two ten-dimensional simplices has larger normalized
projection-body volume than a twenty-dimensional simplex. This gives
a counterexample to Brannen's proposed simplex maximum.
\end{abstract}

\section{The result}\label{sec:result}

\paragraph{Conventions.}
Throughout, \(d\ge2\). A \emph{convex body} in \(\R^d\) is a compact
convex set with nonempty interior. \(|K|\) is the \(d\)-dimensional volume
of \(K\), \(\langle\cdot,\cdot\rangle\) the Euclidean inner product,
\(S^{d-1}\) the unit sphere, \(\proj_{u^\perp}\) the orthogonal projection
onto the hyperplane \(u^\perp\), and \(\vol_{d-1}\) the
\((d-1)\)-dimensional volume. \(B_2^j\) is the Euclidean unit ball of
\(\R^j\) and \(\kappa_j=|B_2^j|\). The \emph{support function} of a compact
convex set \(C\) is \(h_C(x)=\max_{y\in C}\langle x,y\rangle\).

\paragraph{Projection body.}
For a convex body \(K\subset\R^d\), the projection body \(\Pi K\) (written
\(\Pi_dK\) when the dimension needs emphasis) is the compact convex set whose
support function is
\[
 h_{\Pi K}(u)=\vol_{d-1}\bigl(\proj_{u^\perp}K\bigr),
 \qquad u\in S^{d-1},
\]
extended to \(\R^d\) by \(h_{\Pi K}(tu)=t\,h_{\Pi K}(u)\) for \(t\ge0\).
A compact convex set is determined by its support function (fact (S2) in
Section~\ref{sec:tools}). That such a set exists is proved in
Lemma~\ref{lem:facets} for polytopes and in Part~B of
Lemma~\ref{lem:affine} in general. Only polytopes are needed for the main
results.

\paragraph{The functional.}
\begin{equation}\label{eq:functional}
 R_d(K)=\frac{|\Pi K|}{|K|^{d-1}},
 \qquad c_d=\frac{(d+1)d^d}{d!}.
\end{equation}
\(R_d\) is invariant under invertible affine maps (Lemma~\ref{lem:affine}),
and \(c_d\) is its value on every \(d\)-simplex
(Proposition~\ref{prop:simplex}). Brannen conjectured in 1996 that
simplices maximize \(R_d\) among all convex bodies~\cite{Brannen}, that is,
\begin{equation}\label{eq:proposed-bound}
 |\Pi K|\le c_d|K|^{d-1}.
\end{equation}

\begin{theorem}\label{thm:counterexample}
Let \(T_{10}=\conv(0,e_1,\ldots,e_{10})\subset\R^{10}\), where the
\(e_i\) are the coordinate unit vectors, and let
\(K=T_{10}\times T_{10}\subset\R^{20}\). Then
\[
 \frac{R_{20}(K)}{c_{20}}
 =\frac{121\binom{20}{10}}{21\cdot2^{20}}
 =\frac{22\,355\,476}{22\,020\,096}>1.
\]
Thus \(K\) violates \eqref{eq:proposed-bound} in dimension twenty.
\end{theorem}

The ratio is about \(1.0152\). A second result,
Corollary~\ref{cor:exponential-excess} in
Section~\ref{sec:counterexample}, iterates the construction: for some fixed
\(\lambda>1\) and all sufficiently large \(n\) there is a product of
simplices \(K_n\subset\R^n\) with \(R_n(K_n)\ge\lambda^nc_n\).

\paragraph{Context (quoted from the source; not used in any proof).}
In \eqref{eq:proposed-bound} it is the volume of \(\Pi K\) itself that
occurs, not the volume of its polar. The classical projection inequalities
involving the polar concern a different functional. For \(d\ge3\), Petty's
separate minimum problem asks whether ellipsoids are the only minimizers of
the same functional \(R_d\); see \cite[Section~1]{Saroglou} for these
distinctions. In dimension three, Chen, Feng, Li, Xi, and Xu settle the
minimum problem and also prove \(R_3(K)\le18=c_3\) for all
three-dimensional convex bodies, with tetrahedra attaining
equality~\cite[Theorems~1.1 and~1.2]{ChenEtAl}. Feng, Hu, Liu, and Xu give
counterexamples to \eqref{eq:proposed-bound} in every dimension \(d\ge9\),
using polytopes with at most \(d+2\) facets~\cite[Theorem~1.3]{FengEtAl}.
Their construction uses perturbed projections of a cube and Minkowski's
existence theorem~\cite[Section~4]{FengEtAl}. So the failure of the simplex
maximum was already known. The present paper supplies a direct
Cartesian-product construction with an exact value. The optimal upper
constant and its maximizing bodies remain separate questions. For centrally
symmetric bodies, Henk gives the lower bound
\(2^d(9/8)^{\lfloor d/3\rfloor}\) for the optimal upper
constant~\cite{Henk}.

\paragraph{Notation changes from the source.}
To avoid clashes, some symbols are renamed. Facet areas are \(\sigma_F\)
(source: \(s_F\), clashing with the dimension \(s\)). In
Lemma~\ref{lem:affine} the circumradius is \(\rho\) (source: \(R\)), the
projection-volume function is \(g\) (source: \(F\)), and the dilation
factors are \(\mu_m\) (source: \(\lambda_m\)). In
Corollary~\ref{cor:exponential-excess} the residue index is \(q\) (source:
\(r\)) and the constant \(c_{10}/\beta^{10}\) is \(\alpha\) (source: \(A\)).

\section{Proof at a glance}\label{sec:glance}

\paragraph{Key idea.}
The functional \(R\) is \emph{multiplicative} on Cartesian products:
\(R_{r+s}(A\times B)=R_r(A)R_s(B)\). The simplex value \(c_d\) is not
multiplicative in this sense: \(c_{10}^2>c_{20}\). Hence
\(T_{10}\times T_{10}\) beats the \(20\)-simplex.

\paragraph{Steps and dependencies.}
\begin{enumerate}[leftmargin=2em,itemsep=2pt]
\item \textbf{Lemma~\ref{lem:facets} (facet formula).} For a polytope
\(P\), \(\Pi P\) is the Minkowski sum of the segments
\([-\tfrac12\sigma_F\nu_F,\tfrac12\sigma_F\nu_F]\) over the facets \(F\).
This is Cauchy's projection formula. Every later result uses it.
\item \textbf{Proposition~\ref{prop:product} (product identity).} The
facets of \(A\times B\) are \(F\times B\) and \(A\times G\). Feeding this
into Lemma~\ref{lem:facets} gives
\(\Pi(A\times B)=(|B|\Pi A)\times(|A|\Pi B)\), and a volume count gives
multiplicativity.
\item \textbf{Lemma~\ref{lem:affine} (affine invariance).} It follows from
Lemma~\ref{lem:facets} by tracking how facet normals and areas transform.
Its only use is the last sentence of Proposition~\ref{prop:simplex}: every
\(d\)-simplex has the same value of \(R_d\), so that \(c_d\) is ``the
simplex value''. The numerical identity in Theorem~\ref{thm:counterexample}
does not depend on it.
\item \textbf{Proposition~\ref{prop:simplex} (simplex value).} By
Lemma~\ref{lem:facets}, \(\Pi T_d\) is a translate of
\(([0,1]^d+[0,w])/(d-1)!\) with \(w=(1,\dots,1)\). A Fubini argument gives
\(|[0,1]^d+[0,w]|=d+1\), hence \(R_d(T_d)=c_d\).
\item \textbf{Theorem~\ref{thm:counterexample}} is
Proposition~\ref{prop:product} plus Proposition~\ref{prop:simplex} plus
integer arithmetic.
\item \textbf{Corollary~\ref{cor:exponential-excess}} is the same two
propositions applied to \(T_{10}^k\times T_{d}\), \(10\le d\le19\),
together with the elementary bounds \(c_n<(n+1)(11/4)^n\) and
\(c_{10}>(11/4)^{10}\).
\end{enumerate}
\[
\text{Lemma~\ref{lem:facets}}\Longrightarrow
\bigl\{\text{Prop.~\ref{prop:product}},\ \text{Lemma~\ref{lem:affine}},\
\text{Prop.~\ref{prop:simplex}}\bigr\},
\qquad
\text{Prop.~\ref{prop:product}}+\text{Prop.~\ref{prop:simplex}}
\Longrightarrow
\bigl\{\text{Thm.~\ref{thm:counterexample}},\
\text{Cor.~\ref{cor:exponential-excess}}\bigr\}.
\]

\paragraph{One flagged step.}
In the proof of Proposition~\ref{prop:product}, an intermediate displayed
formula of the source has two exponents interchanged. It is marked there
with a box. The proposition itself is unaffected.

\section{Tools: support functions and volumes}\label{sec:tools}

Minkowski addition is \(C+D=\{c+d:c\in C,\ d\in D\}\). The following facts
are used repeatedly; \(C,D\) are nonempty compact convex sets.

\begin{enumerate}[label=\textbf{(S\arabic*)},leftmargin=3em,itemsep=3pt]
\item \emph{Calculus of support functions.} Directly from the definition
(a maximum over independent variables splits):
\(h_{C+D}=h_C+h_D\); \(h_{tC}=t\,h_C\) for \(t\ge0\);
\(h_{MC}(u)=h_C(M^tu)\) for a linear map \(M\); and
\(h_{C\times D}(u,v)=h_C(u)+h_D(v)\) for \(C\subset\R^r\),
\(D\subset\R^s\).
\item \emph{Support functions determine and order sets.} If
\(h_C\le h_D\) everywhere, then \(C\subset D\). In particular
\(h_C=h_D\) implies \(C=D\).

\emph{Proof.} Let \(x\notin D\) and let \(y\in D\) be a point of \(D\)
nearest to \(x\). For \(z\in D\) and \(t\in(0,1]\) the point
\(y+t(z-y)\) lies in \(D\), so
\(\|x-y\|^2\le\|x-y-t(z-y)\|^2
=\|x-y\|^2-2t\langle x-y,z-y\rangle+t^2\|z-y\|^2\).
Dividing by \(t\) and letting \(t\to0\) gives
\(\langle x-y,z-y\rangle\le0\). With \(p=x-y\neq0\), this says
\(\langle p,z\rangle\le\langle p,y\rangle
=\langle p,x\rangle-\|p\|^2<\langle p,x\rangle\) for all \(z\in D\), so
\(h_D(p)<\langle p,x\rangle\). If \(x\) were in \(C\), then
\(\langle p,x\rangle\le h_C(p)\le h_D(p)\), a contradiction. Hence
\(C\subset D\). \(\square\)
\item \emph{Segments.} For \(v\in\R^d\),
\(h_{[-v/2,\,v/2]}(u)=\tfrac12|\langle v,u\rangle|\), because a linear
function on a segment is maximized at an endpoint.
\item \emph{Volumes.} \(|MC|=|\det M|\,|C|\) for linear \(M\);
\(|tC|=t^d|C|\) for \(t\ge0\) and \(C\subset\R^d\);
\(|C\times D|=|C|\,|D|\).
\item \emph{Interior points.} If \(q\) is an interior point of a convex
set \(C\) and \(p\in C\), then \((1-t)p+tq\) is an interior point of \(C\)
for every \(t\in(0,1]\). Indeed, if a ball \(B\) around \(q\) lies in
\(C\), then \((1-t)p+tB\subset C\) is a ball around \((1-t)p+tq\).
\end{enumerate}

\section{Facet formula}\label{sec:facets}

The next lemma writes the projection body of a polytope as a sum of
segments, one per facet. It is the polytopal case of Cauchy's projection
formula; see \cite[Section~1]{Saroglou}. The proof is included because the
factor \(\tfrac12\) affects all later constants. The lemma is used in
Propositions~\ref{prop:product} and~\ref{prop:simplex} and in
Lemma~\ref{lem:affine}.

\begin{lemma}\label{lem:facets}
Let \(P\subset\R^d\), \(d\ge2\), be a full-dimensional convex polytope. For
each facet \(F\), let \(\sigma_F\) be its \((d-1)\)-dimensional volume and
\(\nu_F\) its outward unit normal. Then
\begin{equation}\label{eq:facet}
 h_{\Pi P}(u)=\frac12\sum_F\sigma_F|\langle\nu_F,u\rangle|,
 \qquad u\in\R^d,
\end{equation}
and therefore
\begin{equation}\label{eq:zonotope}
 \Pi P=\sum_F
 \left[-\frac{\sigma_F\nu_F}{2},\frac{\sigma_F\nu_F}{2}\right].
\end{equation}
\end{lemma}

\begin{proof}
\emph{Plan.} For a unit vector \(u\), we show that almost every point of
\(\proj_{u^\perp}P\) is covered exactly twice by the projected facets, and
that the facet \(F\) projects to a set of volume
\(\sigma_F|\langle\nu_F,u\rangle|\). Integrating the covering multiplicity
gives \eqref{eq:facet}. Fix a unit vector \(u\) and write
\(\pi=\proj_{u^\perp}\).

\step{Step 1 (Jacobian).} \emph{Let \(H\) be a linear hyperplane with unit
normal \(\nu\). Then \(\pi|_H\colon H\to u^\perp\) multiplies
\((d-1)\)-volume by \(|\langle\nu,u\rangle|\).}
Let \(v_1,\dots,v_{d-1}\) be an orthonormal basis of \(H\). The factor is
the \((d-1)\)-volume of the parallelepiped spanned by
\(\pi v_1,\dots,\pi v_{d-1}\) in \(u^\perp\). Since \(u\) is a unit vector
orthogonal to all \(\pi v_i\), this equals
\(|\det(\pi v_1,\dots,\pi v_{d-1},u)|\). Because
\(\pi v_i=v_i-\langle v_i,u\rangle u\), column operations turn this into
\(|\det(v_1,\dots,v_{d-1},u)|\). Expanding
\(u=\langle u,\nu\rangle\nu+\sum_i\langle u,v_i\rangle v_i\) in the last
column, the terms with \(v_i\) vanish, and
\(|\det(v_1,\dots,v_{d-1},\nu)|=1\). The result is
\(|\langle u,\nu\rangle|\).

\step{Step 2 (projected facets).} By Step~1 applied to the hyperplane
parallel to that of \(F\),
\begin{equation}\label{eq:projfacet}
 \vol_{d-1}(\pi F)=\sigma_F|\langle\nu_F,u\rangle|
 \qquad\text{for every facet }F.
\end{equation}
If \(\langle\nu_F,u\rangle=0\), the hyperplane of \(F\) contains the
direction \(u\), so \(\pi F\) lies in a \((d-2)\)-dimensional affine
subspace of \(u^\perp\), and both sides are \(0\).

\step{Step 3 (exceptional set).} Let \(E\subset u^\perp\) be the union of
\begin{enumerate}[label=(\roman*),leftmargin=2.5em,itemsep=1pt,topsep=2pt]
\item the boundary of \(\pi P\) relative to \(u^\perp\);
\item the projections of all faces of \(P\) of dimension at most \(d-2\);
\item the projections of all facets \(F\) with
\(\langle\nu_F,u\rangle=0\).
\end{enumerate}
\(E\) has \((d-1)\)-dimensional measure zero. For (ii) and (iii) the sets
lie in finitely many affine subspaces of dimension at most \(d-2\)
(Step~2 for (iii)). For (i), \(C=\pi P\) is a compact convex set with
nonempty interior in \(u^\perp\). After translating an interior point to
\(0\), (S5) gives \((1-\varepsilon)C\subset\operatorname{int}C\), so the
boundary lies in \(C\setminus(1-\varepsilon)C\), of measure
\((1-(1-\varepsilon)^{d-1})\vol_{d-1}(C)\to0\).

\step{Step 4 (multiplicity two).} \emph{Every \(x\in\pi P\setminus E\) lies
in \(\pi F\) for exactly two facets \(F\).}
Let \(\ell=x+\R u\).

(a) \emph{\(\ell\) meets the interior of \(P\).} Pick
\(q\in\operatorname{int}P\). Since \(x\) is interior to \(\pi P\) by (i),
\(x'=x+\delta(x-\pi q)\in\pi P\) for some \(\delta>0\). Choose \(p'\in P\)
with \(\pi p'=x'\). The point \((p'+\delta q)/(1+\delta)\) is interior to
\(P\) by (S5), and it projects to \((x'+\delta\pi q)/(1+\delta)=x\).

(b) Hence \(\ell\cap P\) is a segment \([p_-,p_+]\) with \(p_-\ne p_+\).
By (S5) all points strictly between \(p_-\) and \(p_+\) are interior to
\(P\), while \(p_\pm\in\partial P\).

(c) The boundary of a polytope is the union of its facets, and two distinct
facets meet in a face of dimension at most \(d-2\). By (ii), \(p_+\) lies
in no such face, so \(p_+\) lies in exactly one facet \(F_+\). Likewise
\(p_-\) lies in exactly one facet \(F_-\). By (iii),
\(\langle\nu_{F_\pm},u\rangle\ne0\).

(d) \(F_+\ne F_-\): if \(\langle\nu_F,u\rangle\ne0\), the function
\(\langle\nu_F,\cdot\rangle\) is strictly monotone along \(\ell\), so the
hyperplane of \(F\) meets \(\ell\) in at most one point.

(e) Finally, \(x\in\pi F\) if and only if \(F\) meets \(\ell\). The points
of \(\ell\) in facets are exactly \(\ell\cap\partial P=\{p_-,p_+\}\). So
the facets with \(x\in\pi F\) are exactly \(F_-\) and \(F_+\).

\step{Step 5 (integration).} By Step~4, \(\sum_F\mathbf 1_{\pi F}=2\) on
\(\pi P\setminus E\). The sum is \(0\) outside \(\pi P\), and \(E\) is null
by Step~3. Integrating over \(u^\perp\) and using \eqref{eq:projfacet},
\[
 2\vol_{d-1}(\pi P)=\sum_F\vol_{d-1}(\pi F)
 =\sum_F\sigma_F|\langle\nu_F,u\rangle|.
\]
This is \eqref{eq:facet} for unit \(u\). Both sides of \eqref{eq:facet} are
positively homogeneous of degree one in \(u\) (the left side by
definition), so \eqref{eq:facet} holds for all \(u\in\R^d\).

\step{Step 6 (zonotope).} By (S3) and (S1), the right side of
\eqref{eq:facet} is the support function of the compact convex set on the
right of \eqref{eq:zonotope}. So this set has the support function required
of \(\Pi P\). It therefore \emph{is} \(\Pi P\), and \(\Pi P\) is unique by
(S2).
\end{proof}

\emph{Sanity check.} For the unit square \(Q=[0,1]^2\) and \(u=e_1\):
the projection is a unit segment, of length \(1\). The formula gives
\(\tfrac12(1+1+0+0)=1\).

\section{Cartesian products}\label{sec:product}

The next proposition is the source of the counterexample. It is used in the
proofs of Theorem~\ref{thm:counterexample} and
Corollary~\ref{cor:exponential-excess}.

\begin{proposition}\label{prop:product}
Let \(A\subset\R^r\) and \(B\subset\R^s\) be full-dimensional convex
polytopes, where \(r,s\ge2\). In the orthogonal product space
\(\R^{r+s}=\R^r\times\R^s\),
\begin{equation}\label{eq:product-body}
 \Pi(A\times B)=(|B|\Pi A)\times(|A|\Pi B).
\end{equation}
In particular,
\begin{equation}\label{eq:multiplicative}
 R_{r+s}(A\times B)=R_r(A)R_s(B).
\end{equation}
\end{proposition}

\begin{proof}
\emph{Plan.} Identify the facets of \(A\times B\) with their normals and
areas (Steps 1--3), apply Lemma~\ref{lem:facets} (Steps 4--5), and compare
volumes (Step 6). Below, \(F\) ranges over the facets of \(A\) and \(G\)
over the facets of \(B\).

\step{Step 1 (exposed faces of a product).} For a compact convex set \(C\)
and a vector \(u\), let \(C_u=\{y\in C:\langle u,y\rangle=h_C(u)\}\) be the
face exposed by \(u\). For \((a,b)\in A\times B\),
\(\langle(u,v),(a,b)\rangle=\langle u,a\rangle+\langle v,b\rangle
\le h_A(u)+h_B(v)\), with equality if and only if \(a\in A_u\) and
\(b\in B_v\). Hence
\[
 (A\times B)_{(u,v)}=A_u\times B_v,\qquad
 \dim(A_u\times B_v)=\dim A_u+\dim B_v.
\]

\step{Step 2 (facets).} \(A_0=A\). If \(u\ne0\), then \(A_u\) lies in the
hyperplane \(\{\langle u,\cdot\rangle=h_A(u)\}\), so \(\dim A_u\le r-1\).
The same holds for \(B\). Every facet of a polytope is the face exposed by
its outward unit normal. Let \(\Phi\) be a facet of \(A\times B\) with
outward unit normal \((u,v)\). If both \(u\ne0\) and \(v\ne0\), then by
Step~1 \(\dim\Phi\le(r-1)+(s-1)=r+s-2\), which is impossible for a facet.
So exactly one of \(u,v\) is zero. If \(v=0\), then \(\Phi=A_u\times B\)
with \(\dim A_u=r-1\). So \(A_u\) is a facet \(F\) of \(A\) and
\(u=\nu_F\). Conversely, for each facet \(F\) of \(A\), the set
\(F\times B=(A\times B)_{(\nu_F,0)}\) has dimension \(r+s-1\), so it is a
facet with outward unit normal \((\nu_F,0)\). The case \(u=0\) is
symmetric. Therefore the facets of \(A\times B\) are exactly
\[
 F\times B\ \text{ with normal }(\nu_F,0),
 \qquad
 A\times G\ \text{ with normal }(0,\nu_G).
\]

\step{Step 3 (areas).}
\(\vol_{r+s-1}(F\times B)=\sigma_F|B|\) and
\(\vol_{r+s-1}(A\times G)=|A|\sigma_G\).
Indeed, \(F\times B\) lies in \(H_F\times\R^s\), where \(H_F\) is the
hyperplane of \(F\). Because the two coordinate blocks are orthogonal, an
isometry \(H_F\cong\R^{r-1}\) induces an isometry
\(H_F\times\R^s\cong\R^{r-1}\times\R^s\) carrying \(F\times B\) to a
product set, whose volume is the product of the volumes by Fubini. So the
area-normal vectors (area times outward unit normal) are
\[
 (|B|\sigma_F\nu_F,0)
 \quad\hbox{and}\quad
 (0,|A|\sigma_G\nu_G).
\]

\step{Step 4 (facet formula).} Apply \eqref{eq:facet} to the
full-dimensional polytope \(A\times B\subset\R^{r+s}\), and then to \(A\)
and to \(B\):
\[
 h_{\Pi(A\times B)}(u,v)
 =\frac12\sum_F|B|\sigma_F|\langle\nu_F,u\rangle|
  +\frac12\sum_G|A|\sigma_G|\langle\nu_G,v\rangle|
 =|B|h_{\Pi A}(u)+|A|h_{\Pi B}(v).
\]

\step{Step 5 (identification).} By (S1), the right side is the support
function of \((|B|\Pi A)\times(|A|\Pi B)\) at \((u,v)\). By (S2), this
proves \eqref{eq:product-body}.

\step{Step 6 (volumes).} Since \(|A\times B|=|A||B|\), the denominator of
\(R_{r+s}(A\times B)\) is \((|A||B|)^{r+s-1}\). The source concludes with
the display
\[
 \begin{split}
 R_{r+s}(A\times B)
 &=\frac{|B|^s|A|^r|\Pi A||\Pi B|}
         {(|A||B|)^{r+s-1}}\\
 &=\frac{|\Pi A|}{|A|^{r-1}}
   \frac{|\Pi B|}{|B|^{s-1}}.
 \end{split}
\]
\unverified{The numerator in the first line does not match
\eqref{eq:product-body}. The factor \(|B|\Pi A\) lies in \(\R^r\), so by
(S4) its volume is \(|B|^r|\Pi A|\); likewise \(|A|\Pi B\subset\R^s\) has
volume \(|A|^s|\Pi B|\). The numerator should therefore be
\(|B|^r|A|^s|\Pi A||\Pi B|\), with \(r\) and \(s\) interchanged relative to
the source. As printed, the first line simplifies to
\(|\Pi A||\Pi B|/(|A|^{s-1}|B|^{r-1})\), which equals the second line only
if \(r=s\) or \(|A|=|B|\). This appears to be a transposition slip: the
second line and \eqref{eq:multiplicative} are correct, as the computation
below shows.}
The computation that I can verify is, by \eqref{eq:product-body} and (S4),
\[
 R_{r+s}(A\times B)
 =\frac{|B|^r|\Pi A|\cdot|A|^s|\Pi B|}{(|A||B|)^{r+s-1}}
 =\frac{|\Pi A|}{|A|^{(r+s-1)-s}}\cdot\frac{|\Pi B|}{|B|^{(r+s-1)-r}}
 =\frac{|\Pi A|}{|A|^{r-1}}\cdot\frac{|\Pi B|}{|B|^{s-1}},
\]
which is \eqref{eq:multiplicative}. Theorem~\ref{thm:counterexample} uses
only \(r=s=10\), where the two versions of the numerator coincide.
Corollary~\ref{cor:exponential-excess} uses \(r\ne s\) and relies on
\eqref{eq:multiplicative} itself.
\end{proof}

\section{Affine invariance}\label{sec:affine}

The product identity reduces the comparison to the values of the factors.
To call \(c_d\) ``the simplex value'' we need all \(d\)-simplices to have
the same value of \(R_d\). The affine covariance of \(\Pi\) is due to
Petty~\cite{Petty}; see also \cite[Equation~(1)]{Ludwig}. The lemma is used
only in the last sentence of Proposition~\ref{prop:simplex}, and there only
its polytope case (Part~A) is needed.

\begin{lemma}\label{lem:affine}
Let \(P\subset\R^d\), \(d\ge2\), be a convex body. For every
invertible linear map \(L\colon\R^d\to\R^d\),
\[
 \Pi(LP)=|\det L|L^{-t}\Pi P,\qquad R_d(LP)=R_d(P),
\]
where \(L^{-t}\) denotes the inverse transpose. Translations also
preserve \(R_d\).
\end{lemma}

\begin{proof}
\emph{Plan.} Part~A treats polytopes via Lemma~\ref{lem:facets}: each
area-normal vector \(\sigma_F\nu_F\) is transformed by \(|\det L|L^{-t}\).
Part~B passes to general convex bodies by inscribed polytopes, and also
shows that \(\Pi P\) exists.

\step{Translations.} \(\proj_{u^\perp}(P+x)\) is a translate of
\(\proj_{u^\perp}P\), so \(h_{\Pi(P+x)}=h_{\Pi P}\). Hence
\(\Pi(P+x)=\Pi P\), and \(R_d(P+x)=R_d(P)\).

\medskip\noindent\textbf{Part A: \(P\) a polytope.}

\step{A1 (normals).} Let \(F\) be a facet of \(P\), lying in
\(\{x:\langle\nu_F,x\rangle=t_F\}\), with
\(P\subset\{\langle\nu_F,x\rangle\le t_F\}\). Put \(n_F=L^{-t}\nu_F\).
Then \(\langle n_F,Lx\rangle=\langle\nu_F,x\rangle\), so
\(LP\subset\{\langle n_F,y\rangle\le t_F\}\) and
\(LF=LP\cap\{\langle n_F,y\rangle=t_F\}\). Since \(L\) is a linear
bijection, the facets of \(LP\) are exactly the sets \(LF\), and the
outward unit normal of \(LF\) is \(n_F/\|n_F\|\).

\step{A2 (areas).} \(\sigma_{LF}=|\det L|\,\|n_F\|\,\sigma_F\).
The prism \(Z=F+[0,\nu_F]\) has volume \(\sigma_F\cdot1\), so
\(LZ=LF+[0,L\nu_F]\) has volume \(|\det L|\sigma_F\) by (S4). On the other
hand, \(LZ\) is swept out by translates of \(LF\). By Fubini its volume is
\(\sigma_{LF}\) times its height over the hyperplane of \(LF\), namely
\[
 \Bigl\langle L\nu_F,\frac{n_F}{\|n_F\|}\Bigr\rangle
 =\frac{\langle\nu_F,L^tL^{-t}\nu_F\rangle}{\|n_F\|}
 =\frac1{\|n_F\|}.
\]
Dividing volume by height gives the formula.

\step{A3 (area-normal vectors).} By A1 and A2,
\(\sigma_{LF}\,n_F/\|n_F\|=|\det L|\,L^{-t}(\sigma_F\nu_F)\).
No sign of \(\det L\) enters, so this also holds when \(\det L<0\).

\step{A4 (covariance).} Put \(M=|\det L|L^{-t}\). By \eqref{eq:zonotope}
for \(LP\) and A3,
\(\Pi(LP)=\sum_F[-\tfrac12M\sigma_F\nu_F,\tfrac12M\sigma_F\nu_F]
=M\sum_F[-\tfrac12\sigma_F\nu_F,\tfrac12\sigma_F\nu_F]=M\,\Pi P\).
This uses that a linear map sends segments to segments and commutes with
Minkowski sums.

\step{A5 (invariance).} \(|\det M|=|\det L|^d\,|\det L|^{-1}
=|\det L|^{d-1}\). So \(|\Pi(LP)|=|\det L|^{d-1}|\Pi P|\), while
\(|LP|^{d-1}=|\det L|^{d-1}|P|^{d-1}\). Hence \(R_d(LP)=R_d(P)\).

\medskip\noindent\textbf{Part B: \(P\) a general convex body}
(for completeness; not needed for Theorem~\ref{thm:counterexample} or
Corollary~\ref{cor:exponential-excess}).

\step{B1 (inscribed polytopes).} By translation invariance we may assume
that \(0\) is an interior point of \(P\). Choose \(\rho>0\) with
\(P\subset\rho B_2^d\), and a polytope \(P_0\subset P\) containing a ball
\(aB_2^d\), \(a>0\) (for example a small cube centred at \(0\)). For
\(m\ge1\) let \(\varepsilon_m=1/m\). By compactness choose finite sets
\(N_1\subset N_2\subset\cdots\subset P\) with
\(P\subset N_m+\varepsilon_mB_2^d\), and put \(P_m=\conv(P_0\cup N_m)\).
Then the \(P_m\) are increasing polytopes containing \(aB_2^d\), and
\[
 P_m\subset P\subset P_m+\varepsilon_mB_2^d
 \subset(1+\varepsilon_m/a)P_m,
 \qquad \varepsilon_m\longrightarrow0.
\]
The last inclusion holds because
\(\varepsilon_mB_2^d=(\varepsilon_m/a)\,aB_2^d\subset(\varepsilon_m/a)P_m\)
and \(P_m+tP_m=(1+t)P_m\) for convex \(P_m\) and \(t\ge0\). Put
\(\mu_m=1+\varepsilon_m/a\to1\).

\step{B2 (uniform convergence).} For unit \(u\) write
\(g(u)=\vol_{d-1}(\proj_{u^\perp}P)\) and \(f_m(u)=h_{\Pi P_m}(u)
=\vol_{d-1}(\proj_{u^\perp}P_m)\). Projecting
\(P_m\subset P\subset\mu_mP_m\) and taking \((d-1)\)-volumes gives
\(f_m\le g\le\mu_m^{d-1}f_m\). Since
\(\proj_{u^\perp}P_m\subset\rho B_2^{d-1}\), we have
\(f_m\le\kappa_{d-1}\rho^{d-1}\). Therefore
\[
 0\le g(u)-f_m(u)
 \le(\mu_m^{d-1}-1)f_m(u)
 \le (\mu_m^{d-1}-1)\kappa_{d-1}\rho^{d-1},
\]
and \(f_m\to g\) uniformly on \(S^{d-1}\).

\step{B3 (existence of \(\Pi P\)).} Since \(P_m\subset P_{m+1}\), we have
\(f_m\le f_{m+1}\), so \(\Pi P_m\subset\Pi P_{m+1}\) by (S2). Also
\(\kappa_{d-1}a^{d-1}\le f_m\le\kappa_{d-1}\rho^{d-1}\) on \(S^{d-1}\), so
by (S2) the bodies \(\Pi P_m\) contain \(\kappa_{d-1}a^{d-1}B_2^d\) and lie
in \(\kappa_{d-1}\rho^{d-1}B_2^d\). Let \(Z\) be the closure of
\(\bigcup_m\Pi P_m\). It is convex (an increasing union of convex sets is
convex, and so is its closure), compact, and has nonempty interior. Its
support function is \(\sup_mh_{\Pi P_m}\), which on the sphere equals
\(\lim_mf_m=g\). So \(Z\) is a convex body with the support function
required of \(\Pi P\); that is, \(\Pi P=Z\) exists.

\step{B4 (volume sandwiches).} From \(f_m\le g\le\mu_m^{d-1}f_m\) and (S2),
\[
 \Pi P_m\subset\Pi P\subset\mu_m^{d-1}\Pi P_m,
 \qquad
 |\Pi P_m|\le|\Pi P|\le\mu_m^{d(d-1)}|\Pi P_m|.
\]
Likewise \(|P_m|\le|P|\le\mu_m^d|P_m|\). Consequently
\(\mu_m^{-d(d-1)}R_d(P_m)\le R_d(P)\le\mu_m^{d(d-1)}R_d(P_m)\), so
\(R_d(P_m)\to R_d(P)\).

\step{B5 (covariance).} Applying \(L\) to the sandwich of B1 gives
\(LP_m\subset LP\subset\mu_mLP_m\), with \(0\) interior to \(LP\). B2--B4
used only the following properties of \(P_m\) and \(P\): this sandwich,
that the polytopes \(P_m\) are increasing, that they contain a fixed ball
centred at \(0\), and that \(P\) lies in a fixed ball. All four carry
over. The \(LP_m\) are increasing polytopes because \(L\) preserves
inclusions. They contain \(L(aB_2^d)\supset a'B_2^d\) with
\(a'=a/\|L^{-1}\|\), since \(\|y\|\le a'\) implies \(\|L^{-1}y\|\le a\).
And \(LP\subset\rho'B_2^d\) with \(\rho'=\|L\|\rho\). Here \(\|\cdot\|\)
on maps is the operator norm. So B2--B4 hold with \(LP_m,LP,a',\rho'\) in
place of \(P_m,P,a,\rho\), and with the same \(\mu_m\). In particular
\(h_{\Pi(LP_m)}\to h_{\Pi(LP)}\) pointwise on \(\R^d\) and
\(R_d(LP_m)\to R_d(LP)\). By Part~A and (S1), for every \(u\),
\[
 h_{\Pi(LP_m)}(u)=|\det L|\,h_{\Pi P_m}(L^{-1}u)
 \longrightarrow|\det L|\,h_{\Pi P}(L^{-1}u)
 =h_{|\det L|L^{-t}\Pi P}(u).
\]
The two limits agree, so \(\Pi(LP)=|\det L|L^{-t}\Pi P\) by (S2).

\step{B6 (invariance).} By Part~A, \(R_d(LP_m)=R_d(P_m)\). Letting
\(m\to\infty\) and using B4 for \(P\) and for \(LP\) gives
\(R_d(LP)=R_d(P)\).
\end{proof}

\section{The simplex value}\label{sec:simplex}

We now compute the value to which the product will be compared. The
projection body of a simplex is a cube plus one extra segment. The
proposition is used in Theorem~\ref{thm:counterexample} and
Corollary~\ref{cor:exponential-excess}.

\begin{proposition}\label{prop:simplex}
For \(d\ge2\), the standard simplex
\(T_d=\conv(0,e_1,\ldots,e_d)\subset\R^d\) satisfies
\[
 |T_d|=\frac1{d!},\qquad
 |\Pi T_d|=\frac{d+1}{((d-1)!)^d},\qquad
 R_d(T_d)=c_d.
\]
Every full-dimensional \(d\)-simplex has the same value \(c_d\).
\end{proposition}

\begin{proof}
\emph{Plan.} List the \(d+1\) area-normal vectors of \(T_d\), read off
\(\Pi T_d\) from \eqref{eq:zonotope}, and compute its volume by slicing
along the direction \(w=e_1+\cdots+e_d\).

\step{Step 1 (volume).} \(T_d\) is the pyramid with base
\(T_{d-1}\times\{0\}\) and apex \(e_d\) at height \(1\). Its slice at
height \(t\) is a copy of \((1-t)T_{d-1}\), so
\(|T_d|=\int_0^1(1-t)^{d-1}|T_{d-1}|\,dt=|T_{d-1}|/d\). With
\(|T_1|=1\), induction gives \(|T_d|=1/d!\).

\step{Step 2 (facets).} \(T_d=\{x:x_i\ge0,\ \sum_ix_i\le1\}\) has \(d+1\)
facets.
\begin{itemize}[leftmargin=2em,itemsep=1pt,topsep=2pt]
\item For \(i=1,\dots,d\): the facet \(T_d\cap\{x_i=0\}\) is a copy of
\(T_{d-1}\) in a coordinate hyperplane. It has area \(1/(d-1)!\) and
outward normal \(-e_i\). Its area-normal vector is \(-e_i/(d-1)!\).
\item The facet \(F_0=\conv(e_1,\dots,e_d)\) has outward unit normal
\(w/\sqrt d\). Dropping the last coordinate (projection along \(e_d\))
maps \(F_0\) bijectively onto \(T_{d-1}\), since \(e_d\mapsto0\). By
Step~1 of Lemma~\ref{lem:facets} this projection scales area by
\(|\langle w/\sqrt d,e_d\rangle|=1/\sqrt d\). So
\(\sigma_{F_0}=\sqrt d/(d-1)!\), and the area-normal vector is
\(w/(d-1)!\).
\end{itemize}

\step{Step 3 (the projection body).} By \eqref{eq:zonotope}, with
\(\gamma=1/(d-1)!\),
\[
 \Pi T_d=\sum_{i=1}^d\Bigl[-\frac{\gamma e_i}2,\frac{\gamma e_i}2\Bigr]
 +\Bigl[-\frac{\gamma w}2,\frac{\gamma w}2\Bigr].
\]
The sign of a generator does not matter for a centred segment.
Translating by \(\tfrac\gamma2(e_1+\dots+e_d+w)\) moves each segment to
start at the origin. Since \(\sum_i[0,e_i]=[0,1]^d\), this shows that
\(\Pi T_d\) is a translate of
\[
 \frac{[0,1]^d+[0,w]}{(d-1)!}.
\]

\step{Step 4 (volume of \(Q+[0,w]\), \(Q=[0,1]^d\)).} Let
\(\omega=w/\|w\|\). For a compact convex set \(C\) and \(y\in w^\perp\),
let \(I_C(y)=\{t\in\R:y+t\omega\in C\}\). It is a compact interval,
possibly empty, and by Fubini
\(|C|=\int_{w^\perp}\operatorname{length}(I_C(y))\,dy\). A point
\(y+t\omega\) lies in \(Q+[0,w]\) if and only if
\(y+(t-\tau\|w\|)\omega\in Q\) for some \(\tau\in[0,1]\), that is, if and
only if \(t\in I_Q(y)+[0,\|w\|]\). So if \(I_Q(y)=[a,b]\) then
\(I_{Q+[0,w]}(y)=[a,b+\|w\|]\). The fibre is nonempty exactly when
\(y\in\proj_{w^\perp}Q\), and then it is longer by \(\|w\|\). Hence
\[
 |Q+[0,w]|
 =|Q|+\|w\|\vol_{d-1}(\proj_{w^\perp}Q)
 =1+h_{\Pi Q}(w),
\]
where the last equality is the homogeneous extension in the definition of
\(h_{\Pi Q}\). The cube has, for each \(i\), two facets of area one with
normals \(\pm e_i\), and \(|\langle\pm e_i,w\rangle|=1\). So
\eqref{eq:facet} gives \(h_{\Pi Q}(w)=\tfrac12\cdot2d=d\). Thus
\(|Q+[0,w]|=d+1\). Figure~\ref{fig:extrusion} shows the case \(d=2\).

\step{Step 5 (the value).} By Steps 3--4 and (S4),
\(|\Pi T_d|=(d+1)/((d-1)!)^d\). Dividing by \(|T_d|^{d-1}=(d!)^{-(d-1)}\),
\[
 R_d(T_d)=\frac{(d+1)(d!)^{d-1}}{((d-1)!)^d}
 =\frac{d+1}{d!}\Bigl(\frac{d!}{(d-1)!}\Bigr)^d
 =\frac{(d+1)d^d}{d!}=c_d.
\]

\step{Step 6 (all simplices).} Every full-dimensional \(d\)-simplex is the
image of \(T_d\) under an invertible linear map followed by a translation.
Lemma~\ref{lem:affine} (Part~A and the translation statement) gives the
same value \(c_d\).
\end{proof}

\emph{Sanity check.} For \(d=2\): \(|T_2|=\tfrac12\), \(|\Pi T_2|=3\),
\(R_2(T_2)=6=c_2\). For \(d=3\): \(c_3=4\cdot27/6=18\), the value quoted
from \cite{ChenEtAl} in Section~\ref{sec:result}.

\begin{figure}[tb]
\centering
\begin{tikzpicture}[scale=1.5, line join=round]
 \fill[black!7] (0,0) rectangle (1,1);
 \fill[blue!13] (1,0)--(2,1)--(2,2)--(1,1)--cycle;
 \fill[yellow!23] (0,1)--(1,1)--(2,2)--(1,2)--cycle;
 \draw[thick] (0,0)--(1,0)--(2,1)--(2,2)--(1,2)--(0,1)--cycle;
 \draw[thin] (0,1)--(1,1)--(1,0);
 \draw[thin] (1,1)--(2,2);
 \node at (.5,.5) {\(Q\)};
 \node at (1.5,1) {\(1\)};
 \node at (1,1.5) {\(1\)};
 \node[below left] at (0,0) {\(0\)};
 \node[above right] at (2,2) {\((2,2)\)};
 \draw[->,thick] (2.6,.3)--(3.6,1.3);
 \node[right] at (3.15,.65) {\(w=(1,1)\)};
\end{tikzpicture}
\caption{The planar extrusion \(Q+[0,(1,1)]\), with
\(Q=[0,1]^2\), consists of the unit square and two parallelograms of
area one. Its area is \(3\). The fiber argument in
Proposition~\ref{prop:simplex} proves the corresponding volume \(d+1\)
in every dimension.}
\label{fig:extrusion}
\end{figure}

\section{The counterexample and exponential excess}\label{sec:counterexample}

The geometric work is complete. What remains is arithmetic with the
constants \(c_d\).

\begin{proof}[Proof of Theorem~\ref{thm:counterexample}]
\step{Step 1.} \(T_{10}\) is a full-dimensional convex polytope in
\(\R^{10}\), so \(K=T_{10}\times T_{10}\) is compact and convex with
nonempty interior in \(\R^{20}\). Proposition~\ref{prop:product} applies
with \(r=s=10\), and with Proposition~\ref{prop:simplex} gives
\[
 R_{20}(K)=R_{10}(T_{10})^2=c_{10}^2
         =\left(\frac{11\cdot10^{10}}{10!}\right)^2.
\]

\step{Step 2.} Divide by \(c_{20}=21\cdot20^{20}/20!\), and use
\(10^{20}/20^{20}=2^{-20}\) and \(20!/(10!)^2=\binom{20}{10}\):
\[
 \frac{R_{20}(K)}{c_{20}}
 =\frac{121\cdot10^{20}\cdot20!}
        {(10!)^2\cdot21\cdot20^{20}}
 =\frac{121\binom{20}{10}}{21\cdot2^{20}}.
\]

\step{Step 3.} \(\binom{20}{10}=184\,756\) and \(2^{20}=1\,048\,576\), so
the numerator is \(121\cdot184\,756=22\,355\,476\) and the denominator is
\(21\cdot1\,048\,576=22\,020\,096\). The exact integer certificate is
\[
 121\cdot184\,756-21\cdot1\,048\,576=335\,380>0.
\]
So \(R_{20}(K)>c_{20}\), that is, \(|\Pi K|>c_{20}|K|^{19}\). This is a
strict violation of \eqref{eq:proposed-bound}.
\end{proof}

\paragraph{The general product test.}
For \(r,s\ge2\), the same two propositions give
\(R_{r+s}(T_r\times T_s)=c_rc_s\). Dividing by \(c_{r+s}\) and collecting
\(\frac{(r+s)!}{r!\,s!}=\binom{r+s}{r}\),
\[
 \frac{R_{r+s}(T_r\times T_s)}{c_{r+s}}
 =\frac{(r+1)r^r}{r!}\cdot\frac{(s+1)s^s}{s!}\cdot
  \frac{(r+s)!}{(r+s+1)(r+s)^{r+s}}
 =\frac{(r+1)(s+1)}{r+s+1}
   \binom{r+s}{r}\frac{r^rs^s}{(r+s)^{r+s}}.
\]
Thus any universal upper constant for \(R_{r+s}\) must be at least
\(c_rc_s\). Theorem~\ref{thm:counterexample} is the case \(r=s=10\).

\medskip
Iterating the product identity gives the following quantitative
consequence.

\begin{corollary}\label{cor:exponential-excess}
There exist a constant \(\lambda>1\) and an integer \(N\ge20\) such that,
for every integer \(n\ge N\), some full-dimensional convex polytope
\(K_n\subset\R^n\) satisfies
\[
 \frac{R_n(K_n)}{R_n(T_n)}\ge\lambda^n.
\]
The polytopes \(K_n\) may be chosen as Cartesian products of simplices.
\end{corollary}

\begin{proof}
\emph{Plan.} Take \(K_n\) to be a product of many copies of \(T_{10}\) and
one simplex of dimension between \(10\) and \(19\). Its value is a product
of simplex constants. We compare it with \(c_n\) using the upper bound
\(c_n<(n+1)\beta^n\) and the lower bound \(c_{10}>\beta^{10}\), where
\(\beta=11/4\).

\step{Step 1 (the bodies).} For \(q\in\{0,\ldots,9\}\) put \(d_q=10+q\).
Every integer \(n\ge20\) has a unique representation \(n=10k+d_q\) with
\(k\ge1\): \(q\) is the last digit of \(n\) and
\(k=\lfloor n/10\rfloor-1\). In the orthogonal coordinate product space,
take
\[
 K_n=T_{10}^{\,k}\times T_{d_q},
\]
where \(T_{10}^{\,k}\) is the \(k\)-fold Cartesian product.

\step{Step 2 (its value).} A product of full-dimensional convex polytopes
is a full-dimensional convex polytope, and all dimensions involved are at
least \(10\ge2\). Induction on \(k\) with Proposition~\ref{prop:product}
(applied to \(T_{10}^{\,j}\times T_{10}\), then to
\(T_{10}^{\,k}\times T_{d_q}\)) and Proposition~\ref{prop:simplex} give
\[
 R_n(K_n)=c_{10}^k c_{d_q}.
\]

\step{Step 3 (\(e<\beta=11/4\)).} In the series
\(e=\sum_{j\ge0}1/j!\), the terms with \(j\le2\) sum to \(\tfrac52\). For
\(j\ge3\), the ratio of consecutive terms is \(1/(j+1)\le\tfrac14<\tfrac13\),
so \(1/j!\le\tfrac16\,3^{-(j-3)}\), with strict inequality for \(j\ge4\).
Hence
\[
 e=\frac52+\sum_{j=3}^{\infty}\frac1{j!}
 <\frac52+\frac16\sum_{j=0}^{\infty}3^{-j}
 =\frac52+\frac16\cdot\frac32
 =\frac{11}4=\beta.
\]

\step{Step 4 (upper bound for \(c_n\)).} The positive term \(n^n/n!\) of
the series \(e^n=\sum_jn^j/j!\) is strictly smaller than the full sum.
With Step~3,
\[
 c_n=(n+1)\frac{n^n}{n!}<(n+1)e^n<(n+1)\beta^n\qquad(n\ge2).
\]

\step{Step 5 (the ten-dimensional block).} Using
\(10!=2^8\cdot3^4\cdot5^2\cdot7\) and \(10^{10}=2^{10}5^{10}\),
\[
 c_{10}=\frac{11\cdot10^{10}}{10!}
 =\frac{11\cdot2^2\cdot5^8}{3^4\cdot7}
 =\frac{17\,187\,500}{567}.
\]
With \(\beta^{10}=11^{10}/2^{20}\), \(11^9=2\,357\,947\,691\) and
\(2^{22}5^8=2^{14}\cdot10^8\),
\[
 \alpha:=\frac{c_{10}}{\beta^{10}}
 =\frac{2^{22}\cdot5^8}{567\cdot11^9}
 =\frac{1\,638\,400\,000\,000}{1\,336\,956\,340\,797}>1 .
\]
Numerically, \(\alpha\approx1.2255\).

\step{Step 6 (choice of \(\lambda\)).} Fix \(\lambda\) with
\(1<\lambda<\alpha^{1/10}\), and put \(\theta=\alpha/\lambda^{10}>1\).

\step{Step 7 (comparison).} By Proposition~\ref{prop:simplex},
\(R_n(T_n)=c_n\). By Steps 2 and 4, and writing \(n=10k+d_q\) so that
\(\beta^n\lambda^n=(\beta\lambda)^{10k}(\beta\lambda)^{d_q}\) and
\(c_{10}^k/(\beta\lambda)^{10k}=\theta^k\),
\[
 \frac{R_n(K_n)}{R_n(T_n)\lambda^n}
 >\frac{c_{10}^k\,c_{d_q}}{(n+1)\beta^n\lambda^n}
 =c_{d_q}\beta^{-d_q}\lambda^{-d_q}\,
   \frac{\theta^k}{10k+d_q+1}.
\]

\step{Step 8 (thresholds).} For each fixed \(q\), the right side tends to
infinity as \(k\to\infty\), since \(\theta>1\) and exponential growth
dominates the linear denominator. So there is \(k_q\ge1\) such that the
right side is at least \(1\) for all \(k\ge k_q\). Let
\(N=\max_{0\le q\le9}(10k_q+d_q)\ge20\). If \(n\ge N\) and
\(n=10k+d_q\), then \(10k+d_q\ge10k_q+d_q\), so \(k\ge k_q\). Step~7 then
gives \(R_n(K_n)/R_n(T_n)>\lambda^n\).
\end{proof}

\section{Verification checklist}\label{sec:checklist}

The load-bearing steps, in the order a sceptical reader should check them:

\begin{enumerate}[leftmargin=2em,itemsep=3pt]
\item \textbf{The factor \(\tfrac12\) in the facet formula.}
Lemma~\ref{lem:facets}, Steps 3--5: almost every point of the projection
is covered by exactly two projected facets. All constants depend on this.
\item \textbf{Facets of a product and the exponent bookkeeping.}
Proposition~\ref{prop:product}, Steps 2--3 (the facets are \(F\times B\)
and \(A\times G\), with areas \(\sigma_F|B|\) and \(|A|\sigma_G\)) and
Step~6. Step~6 contains the one boxed discrepancy with the source: its
numerator \(|B|^s|A|^r\) should be \(|B|^r|A|^s\). Check that
\((|B|\Pi A)\subset\R^r\) has volume \(|B|^r|\Pi A|\) and that
\eqref{eq:multiplicative} follows.
\item \textbf{The simplex value.} Proposition~\ref{prop:simplex}: Step~2
(the slanted facet has area-normal vector \(w/(d-1)!\)), Step~4
(\(|[0,1]^d+[0,w]|=d+1\)), and the algebra in Step~5. The cases \(d=2,3\)
give \(6\) and \(18\).
\item \textbf{The integer inequality.} Proof of
Theorem~\ref{thm:counterexample}, Steps 2--3:
\(c_{10}^2/c_{20}=121\binom{20}{10}/(21\cdot2^{20})\) and
\(121\cdot184\,756-21\cdot1\,048\,576=335\,380>0\).
\item \textbf{The corollary's two numerical inputs.}
Corollary~\ref{cor:exponential-excess}, Step~3 (\(e<11/4\)) and Step~5
(\(\alpha=2^{22}5^8/(567\cdot11^9)>1\)), and the exponent bookkeeping in
Step~7.
\item \textbf{Role of affine invariance.} Lemma~\ref{lem:affine} is needed
only to say that \emph{every} simplex has value \(c_d\)
(Proposition~\ref{prop:simplex}, Step~6), and for that Part~A suffices.
Part~B (general convex bodies) is not used by the main results.
\end{enumerate}

\begin{thebibliography}{7}

\bibitem{Brannen}
N.~S.~Brannen,
\emph{Volumes of projection bodies},
Mathematika \textbf{43} (1996), no.~2, 255--264.
\href{https://doi.org/10.1112/S002557930001175X}{doi:10.1112/S002557930001175X}.

\bibitem{ChenEtAl}
S.~Chen, Y.~Feng, Y.~Li, D.~Xi, and L.~Xu,
\emph{Petty's conjectured projection inequality in dimension three},
preprint, 2026, MathSciDoc:2608.23001.
\href{https://archive.ymsc.tsinghua.edu.cn/pacm_download/741/12780-pettyinequality.pdf}{Full text}.

\bibitem{FengEtAl}
Y.~Feng, S.~Hu, W.~Liu, and L.~Xu,
\emph{On the reverse projection inequality},
preprint, 2026, MathSciDoc:2608.23002.
\href{https://archive.ymsc.tsinghua.edu.cn/pacm_download/743/12781--2026.8.26.pdf}{Full text}.

\bibitem{Henk}
M.~Henk,
\emph{Note on projection bodies of zonotopes with \(n+1\) generators},
Acta Mathematica Scientia \textbf{45} (2025), 96--103.
\href{https://doi.org/10.1007/s10473-025-0107-9}{doi:10.1007/s10473-025-0107-9}.

\bibitem{Ludwig}
M.~Ludwig,
\emph{Projection bodies and valuations},
Advances in Mathematics \textbf{172} (2002), no.~2, 158--168.
\href{https://doi.org/10.1016/S0001-8708(02)00021-X}{doi:10.1016/S0001-8708(02)00021-X}.

\bibitem{Petty}
C.~M.~Petty,
\emph{Projection bodies},
in \emph{Proceedings of the Colloquium on Convexity (Copenhagen, 1965)},
K\o benhavns Universitet Matematiske Institut, Copenhagen, 1967,
pp.~234--241.

\bibitem{Saroglou}
C.~Saroglou,
\emph{On the shape of a convex body with respect to its second projection body},
Advances in Applied Mathematics \textbf{67} (2015), 55--74.
\href{https://doi.org/10.1016/j.aam.2015.03.004}{doi:10.1016/j.aam.2015.03.004}.
Preprint:
\href{https://arxiv.org/abs/1409.4347v2}{arXiv:1409.4347v2}.

\end{thebibliography}

\end{document}
